\begin{frame}[t]{Un(e) candidat(e) idéal(e)}
  \footnotesize
  \begin{columns}
    \column{0.5\textwidth}
      Vous avez de la chance!
      Une élection hypothétique arrive, et votre candidat.e idéal.e est sur le bulletin de vote.
      Comment est ce.tte candidat.e?
      Réfléchissez aux sujets suivants, puis écrivez un paragraphe pour décrire votre candidat.e idéal.e.
    \column{0.5\textwidth}
      \begin{center}
        \includegraphics[scale=0.08]{candidat_ideal.jpg}
      \end{center}
  \end{columns}
  \only<1>{
    \begin{columns}[t]
      \column{0.5\textwidth}
        \begin{itemize}
          \item Est-ce que le/la candidat.e est d'un parti spécifique?
          \item Le/La candidat.e est de droite, de gauche ou centriste?
          \item Le/La candidat.e s'intéresse à quels problèmes? Quel est son programme?
          \item Est-ce que le/la candidat.e est riche ou pauvre?
        \end{itemize}
      \column{0.5\textwidth}
        \begin{itemize}
          \item Est-ce que le/la candidat.e a une identité spécifique? Un genre? Une race ou ethnicité?
          \item Est-ce que la religion est importante? La laïcité?
          \item Est-ce que c'est un.e citoyen.ne typique ou un.e politicien.ne de carrière?
        \end{itemize}
    \end{columns}
  }
  \only<2>{
    \vspace{1cm}
    Maintenant, discutez de vos candidat.es idéal.es avec des partenaires.
    Est-ce que vous votez pour des candidats similaires ou différentes? Pourquoi?
    \emph{Après finir}, discutez de vos candidat.es avec encore une autre personne de la même façon, puis avec une autre personne, etc.
  }
\end{frame}